\section{模型检验与灵敏度分析}
\subsection{硬约束、复算与双运行}
solve v12 的约束审计显示：任务唯一性、到达、截止、实时开工、时延、GPU、IT 与设施功率均通过；Q3 和 Q4 的能源平衡、SOC、功率、互斥及终端状态最大绝对残差不超过 $2.274\times10^{-13}$。这些检验说明结果在给定容差与题面参数内自洽，但不提供最优性证书。

两次隔离运行使用相同输入副本与同一 code v15。冻结程序的完整 SHA-256 为 \texttt{846f527b60b8365c\allowbreak{}619b806e7171f240\allowbreak{}d010134639386134\allowbreak{}808c2323007a34be}；自然耗时分别为 1687.58 s 和 1706.49 s。除 \texttt{method\_runtime.json} 的实际耗时外，48 个正式输出文件逐字节一致；原始 \texttt{results.json} 的 SHA-256 为 \texttt{691e95d6f0450c63\allowbreak{}c6b0468a55453af7\allowbreak{}78dfe949316eadad\allowbreak{}c63a5ad8a0a4daa0}。这证明给定环境与固定种子下结果可复现。

\subsection{问题二参数灵敏度}
以平衡方案成本 $-421801727.45$ 元为基准，对购电价格整体缩放。系数为 0.8、0.9、1.1、1.2 时，成本分别为 $-421804612.25$、$-421803169.85$、$-421800285.05$、$-421798842.65$ 元，相对变化绝对值不超过 0.000684\%。图~\ref{fig:price} 表明当前净成本对购电价格整体缩放不敏感，原因是购电支出相对售电收入很小。
\begin{figure}[H]\centering
\includegraphics[width=.80\textwidth]{sensitivity_electricity_price.png}
\caption{问题二净成本的购电价格灵敏度}\label{fig:price}
\end{figure}

新能源水平系数为 0.8、0.9、1.1、1.2 时，成本分别为 $-303912263.19$、$-381129183.96$、$-446460089.89$、$-457667554.37$ 元，相对基准变化为 27.9490\%、9.6426\%、$-5.8460$\%、$-8.5030$\%。图~\ref{fig:renewable} 呈现明显非线性：新能源下降时购电缺口扩大，而上升时外送上限与弃电共同限制边际收益。
\begin{figure}[H]\centering
\includegraphics[width=.80\textwidth]{sensitivity_renewable_level.png}
\caption{问题二净成本的新能源水平灵敏度}\label{fig:renewable}
\end{figure}

\subsection{终端 SOC 与高循环策略解释}
Q3 成本最小策略在终端 SOC 等于初值时，成本为 $-458752156.98$ 元、碳排放为 0、等效循环合计 240.186；放宽终端 SOC 后，其成本仍为 $-458752156.98$ 元。由此可见，当前成本最小候选的主要驱动是反复利用新能源与价格结构，而非末端库存约束。平衡与新能源策略的等效循环分别达到 1568.242 和 1568.882，明显高于短期物理运行中通常期望的水平。由于题面未提供循环寿命与退化成本，程序没有据此否决；论文将其列为工程外推的首要限制，而不把高循环账面收益直接解释为可部署收益。

Q4 的三组新能源波动种子产生不同峰值和碳排放：种子 2026、2027、2028 的峰值购电分别为 273.9082、298.6499、146.1294 MW，碳排放分别为 760.7500、852.0600、464.9554 tCO$_2$。这说明单一随机轨迹不足以概括系统风险，当前三种子结果只提供方向性稳健性证据。
